% ===== §16 Reproduction: Repetition without Return =====
\section{Repetition without Return}
\label{sec:reproduction}

\epigraph{\zh{苟日新，日日新，又日新。}}{\textit{The Great Learning}}

\noindent
The fourth movement completes the cycle and is the one most liable to be misread, because the word that names it carries two opposite meanings and the whole account depends on holding them apart. Reproduction can mean the making-again of the same, the duplication that runs an identical operation a second time, and in that meaning it is the mark of the mechanical, the assembly line, the labour of the one condemned to roll the same stone forever. Reproduction can also mean the thickening by which what was shared settles as a raised starting point, so that the next perception begins higher and finer, and in that meaning it is the very accumulation that distinguishes a good cycle from a wheel. The fourth step is reproduction in the second sense and not the first, and the labour of this section is to say, with precision, what divides them, since they can look identical from outside and the difference between them is the difference between a life that gains and a life that merely turns. This is where the metaphysics of repetition, which the earlier treatment of maintenance opened, is brought to its formal point, and where the political economy of the reproductive act, which the next part will take up, is first joined to it.

\subsection{The semantic pivot: reproduction returns to the root, not the origin}
\label{sub:reproduction-pivot}

The distinction the section needs is the one this series has carried as its central geometric figure. To reproduce, in the sense the model requires, is not to repeat the identical but to return to the root without returning to the origin. The cycle closes; perception, reception, sharing, and reproduction bring the practitioner back to perception. But because the sharing has settled as a changed disposition of the field, what is returned to is not the point of departure. The threshold of joint perception has been lowered; the partners can now register what before went unseen; the starting point of the next cycle has been raised. The closing-without-returning is the anholonomy the series has named throughout, the gain that a cycle accumulates precisely by not coming back to where it began, and reproduction is the name, at this most intimate scale, for the settling of that gain.

\claim{Reproduction returns to the root and not to the origin. The good cycle closes, bringing the practitioner back to perception, but the starting point has been raised by what the previous cycle settled, so that what is returned to is the root and not the origin. Reproduction in the required sense is this raising of the starting point, the anholonomy of the intimate cycle; reproduction in the degenerate sense is the closing of the cycle upon its own origin, the wheel whose holonomy is zero. The two are indistinguishable in any single act and are told apart only by whether the next cycle begins higher.}

This pivot fixes the criterion the section will use throughout, and it is a criterion in the dynamical sense and not the evaluative one. Whether a given repetition is good or degenerate is not read from how the act looks, since the same gesture serves both, but from whether accumulation occurs across cycles, whether the spiral gains altitude or the wheel turns in place. The criterion is the one the series proposed for cycles in general, here brought to the scale of a single returned gesture between two people, and the traditions gathered below are gathered to deepen it, to say in their several languages what makes a repetition gain rather than turn.

\begin{casebox}[The same walk, on the hundredth evening]
The two have a walk they take in the evening, the same streets, the same turning, the same return. On the hundredth evening it is, to all appearance, the walk it was on the first: the same route, the same length, the same hour. Yet two hundredth evenings are possible, and they are not the same evening. On the one, the walk has become a track worn smooth, a motion performed because it is performed, each evening a copy of the last, the streets no longer seen, the company no longer met, the walk a wheel that turns because it has always turned. On the other, the hundredth walk is thick with the ninety-nine, the same turning met with a perception deepened by every previous turning, a small new thing seen because so much has been seen before, the company known more finely for the hundred evenings of knowing, so that the walk gains, the field it builds accumulating across the evenings rather than flattening. The route is identical; the repetition is opposite. What divides the deadening walk from the deepening one is the whole question of reproduction, and it cannot be read from the route, which is the same, but only from whether the hundredth evening begins higher than the first.
\end{casebox}

\subsection{The two repetitions}
\label{sub:reproduction-two-repetitions}

The first body of theory is the philosophical analysis of repetition itself, and it converges, from two directions, on the very distinction the pivot drew. One line opposes a backward repetition, the recollection that retrieves what already was, to a forward repetition that produces the new in the form of the again, a repetition by which what is repeated is not retrieved but renewed and gained~\citep{kierkegaard_repetition}, taken up into a future rather than recovered from a past. Reproduction in the model's sense is forward in just this way: the shared beauty is not retrieved unchanged but taken up so that the relation gains, and the gain is what makes the again a thickening rather than a recovery. A second line distinguishes a bare repetition of the identical, which is the mechanical and the dead, from a repetition that always carries difference, in which what returns returns differently and the difference is productive~\citep{deleuze_difference}; on this analysis the mechanical same is the degenerate limit of a repetition whose truth is the carrying of difference. Both lines give philosophical standing to the distinction between the spiral and the wheel, and both may be read alongside the holonomy criterion without being identified with it, since a literary or metaphysical analysis of repetition and a dynamical measure of path-dependent gain are not the same thing and are not made the same by agreeing.

What these analyses contribute is the demonstration that the good repetition is not a special achievement laid over ordinary repetition but the truth of repetition itself, of which the mechanical same is the falling-away. This matters to the paper's politics as much as its metaphysics, because it means that the bad repetition, the wheel, is not the natural state of the everyday from which beauty must rescue it, but a degeneration of the everyday's own generative tendency, a degeneration with causes that the next account locates. The everyday does not have to be redeemed from repetition. It has to be kept from the degeneration of its repetition, and the difference between those two tasks is the difference between an aesthetic education that despises ordinary life and one that recovers it.

\subsection{The political economy of the reproductive act}
\label{sub:reproduction-political-economy}

The word reproduction does not come without its history, and the section would be evading its most serious question if it did not hear, in the term, the analysis of social and laboring reproduction that the term carries. That analysis asks under what conditions the making-again of life becomes the alienated making-again of the same, the worker's gesture repeated for another's accumulation, the stone rolled for a purpose that is never the roller's own. Set beside the intimate cycle, the question becomes sharp and unavoidable: what distinguishes the reproductive act that alienates, the gesture endlessly repeated so that holonomy is siphoned off to a power outside the relation, from the reproductive act that accumulates, the gesture finely adjusted again and again so that the holonomy settles into the relation itself. The two are formally alike. Both are repetitions, both are labour, both make again. They differ in the destination of what the labour produces, in whose field gains the curvature, in whether the accumulation returns to those whose coupling generated it or is captured by an order that owns the channel of its production.

\claim{What makes a reproductive act alienated is not repetition but the capture of its holonomy by an external power. A gesture repeated for another's accumulation and a gesture finely adjusted for the relation's own thickening are formally alike as repetitions; they differ in whether the accumulated value returns to the coupling that generated it or is siphoned to an order outside it. Alienation is therefore not a property of the everyday's repetitiousness but of the ownership of its yield, and the same act is bad repetition or good according to where its holonomy goes.}

This is the section's guard against its own worst temptation, the temptation to preach a serene aesthetics of the dishcloth to people whose repetition is not theirs to shape. It locates the difference between good and bad repetition not in the inwardness of the laborer, as though the worn-down had only to find beauty in their toil, but in the structure of ownership that decides where the holonomy of their labour goes. The cause this account names, the capture of the cycle's yield by an external power, is the same cause the series identified as the most insidious of the external sources of estrangement, the bleeding of a possessive order into a generative one, and it will be carried directly into the political economy of the field, where the question of who has the leisure and the standing to shape their own repetition is met without evasion.

\subsection{An anthropological voice: ritual and the cultivation of repetition}
\label{sub:reproduction-ritual}

The study of ritual offers a third account of how repetition is made to gain rather than dull, and it is worth admitting because it describes, at the scale of cultures, a technology for exactly the good repetition the section theorises. Ritual is repetition deliberately cultivated: the same acts performed again on the recurring occasion, the same words said, the same gestures made, and yet the developed traditions of ritual are emphatic that the repetition is not to be mechanical, that the rite performed by rote is a dead rite, and that the living rite is the one entered freshly each time, the same form filled with a present attention that makes the hundredth performance more than the first. The anthropological literature describes the means by which traditions guard their rituals against the slide into rote, the framing that marks the ritual time as set apart and so demands a heightened attention, the small variations that keep the form from becoming wholly predictable, the investment of the repeated act with a significance that rewards fresh attention. These are cultural technologies of good repetition, devices by which a community keeps its repeated forms gaining rather than dulling, and they model, at the scale of the culture, what the beauty of maintenance performs at the scale of the couple. The shared evening walk, the daily greeting, the repeated meal are the couple's rituals, and the means by which they are kept fresh, the attention that enters them as set apart, the small variations that keep them from rote, the significance with which they are invested, are the dyadic versions of the cultural technologies the study of ritual describes.

The voice contributes a specific insight the other accounts lack, that good repetition is often not spontaneous but cultivated by form, that the freshness of the repeated act is frequently secured precisely by a structure deliberately maintained, the ritual frame that demands attention rather than leaving it to arise. This corrects a possible misreading of the beauty of maintenance, the misreading on which good repetition would be a matter of spontaneous feeling, the partners simply happening to bring fresh attention to the repeated act. The anthropology of ritual suggests otherwise, that the freshness of repetition is often the achievement of a form that secures it, a structure the partners maintain precisely in order that their repetition shall gain, and that the cultivation of maintenance may therefore include the deliberate making of small rituals, forms set apart and invested with significance, whose function is exactly to keep the repeated act from dulling. Good repetition, on this account, is not only a spontaneous freshness but sometimes a cultivated one, secured by forms the partners build for the purpose, and the beauty of maintenance includes the craft of building such forms, the small domestic rituals whose office is to make the daily return the gaining kind.

\subsection{Sedimentation and renewal: an Eastern principal voice and a phenomenological one}
\label{sub:reproduction-sedimentation}

The last body of theory concerns how reproduction accumulates in time without becoming a hoarded stock, and here an Eastern voice and a phenomenological one speak to the same point. The phenomenological account describes how the passing present is retained and carried forward, how experience sediments into a sedimented ground that is not a store of objects but a changed disposition out of which the next experience arises. The accumulation of the good cycle is sedimentation in this sense: not a treasury of kept beauties but a settled ground that raises the next beginning, which is exactly the changed curvature the reception section distinguished from a possessive stock. The Eastern voice gives the same accumulation its ethical accent. The injunction of daily renewal, if truly renewed for a day, then day after day renewed~\citep{daxue}, and again renewed, names a repetition that is not the recurrence of the same but the unceasing freshening of the same matter, the same ordinary act made new each time. Daily renewal is good repetition stated as a discipline of life: not the seeking of ever-new objects, which is the exhaustion the next part diagnoses as the consumption of beauty, but the renewing of the same daily thing so that it thickens rather than dulls.

The provenance of the injunction sharpens its meaning. The words are recorded as an inscription on a washing vessel, a basin used daily for an act as ordinary and as repetitive as any in a life, and the placing of the maxim there is itself the teaching: the renewal it commands is to be found in the most repeated and least exalted of daily acts, the daily washing, and not in some exceptional occasion of renewal set apart from ordinary life. To inscribe make it new, day by day, on the basin one uses every morning is to say that the newness is the newness of the same daily thing, the same washing renewed rather than replaced, and that the site of renewal is exactly the repetitive everyday the tradition might have been expected to despise. This is the doctrine of good repetition fixed to its humblest possible object, and it confirms from the Eastern side the whole orientation of this paper: that the renewal which keeps a life from dulling is not the escape from the daily and repetitive into the novel and exceptional, but the freshening of the daily and repetitive itself, the basin made new each morning by the manner of its use rather than by its replacement with a finer basin. The walk on the hundredth evening is renewed in just this sense, made new not by a new route but by the freshening of the same route, which is the only renewal a repetitive life affords and, the inscription insists, the only one it needs.

The two voices together resolve a worry the reception section deliberately left standing, the worry that accumulation must mean storing and so must violate the mirror that does not abide. Sedimentation is accumulation that does not store, because what settles is not a stock in a place but the raised ground of the next arising; daily renewal is repetition that does not dull, because what returns is freshened rather than recurred. The accumulation reproduction requires is therefore exactly the non-possessive, non-dulling accumulation these voices describe, the settling of curvature and the freshening of the same, and this is why responding-yet-not-storing, far from forbidding reproduction, was already its condition. The mirror was a teaching about time, and these voices say what the teaching, carried into time, becomes: a ground that rises without being hoarded and a sameness that is made new without being merely repeated.

\subsection{Two real conflicts}
\label{sub:reproduction-conflicts}

The accounts of reproduction conflict at two points, and the second closes the figure that has run through all four sections.

The first conflict reprises, at the level of time, the tension the reception section met: accumulation against non-abiding. The philosophical and phenomenological accounts speak of thickening, sedimentation, a raised ground, a gain carried forward; the contemplative voice that governed reception spoke of a mirror that retains nothing. Stated baldly the two seem to demand opposite things, that reproduction lay up and that nothing be laid up. The resolution was prepared in the reception section and is completed here: what the mirror forbids is the possessive stock, the beauty seized and kept in a place against loss, and what reproduction accumulates is not a stock but a changed curvature, a ground that has risen, a threshold that has lowered. These are not the same kind of thing, and only the first is storing. Sedimentation and daily renewal name precisely the accumulation that is not a hoard, the settling of disposition and the freshening of the same, so that the conflict between accumulation and non-abiding dissolves once it is seen that the accumulation in question was never of the kind non-abiding forbids. The mirror and the rising ground are consistent because one speaks of what must not be grasped and the other of what cannot be grasped, the curvature of a path being no candidate for a hoard.

The second conflict is the gravitational one, the pull of the wheel. Even granting the distinction between good and degenerate repetition, the degenerate is not a remote possibility but a constant tendency; repetition slides toward the mechanical, the shared dulls into routine, the spiral loses altitude and flattens toward the wheel. If this drift is the natural fate of repetition, then the good cycle is at best a temporary victory against an entropy that will reclaim it, and reproduction, the step that was to secure accumulation, is forever undermined by the very repetitiousness it requires. The conflict is not idle. Anyone who has watched a shared life dull into habit knows the pull is real, and an account that merely asserted the good repetition without facing the drift toward the bad would be sentimental in exactly the way the paper has tried not to be.

\subsection{The hardest step: reproduction is fed by the prior three and so closes the cycle}
\label{sub:reproduction-hardest}

The resolution to the second conflict is the one that completes the model, and it turns on seeing that reproduction is not a step performed alone. The drift toward the wheel is real for any repetition considered in isolation, a making-again with no fresh input, and such a repetition does flatten, because nothing enters it to renew what it repeats. But reproduction in the model is not an isolated making-again. It is the fourth movement of a cycle whose first three movements feed it: each turn of the spiral is renewed by a fresh perception that sees what before went unseen, a reception that lets the new alteration in, a sharing that generates between the two a beauty that did not exist before. The thickening of reproduction is not the step's own solitary achievement, a stock it builds by repeating; it is the settling of what the prior three steps have freshly produced this turn. Reproduction flattens toward the wheel exactly when the prior three have ceased, when nothing new is perceived, received, or shared, and the making-again runs on empty. It gains altitude exactly when they continue, when each turn carries new contingency in.

\claim{Reproduction does not secure accumulation on its own and cannot. It thickens only insofar as the prior three movements continue to feed it, a fresh perception, a new reception, a generated sharing, on each turn; deprived of them it runs on empty and flattens toward the wheel. The health of reproduction is therefore not a property of the fourth step but a function of whether the first three are still occurring. This is what closes the four movements into a cycle rather than a line: reproduction returns not to a stored deposit but to perception, and the return is alive only if perception is alive. The drift toward the wheel is resisted not by working harder at reproduction but by keeping the cycle turning.}

An objection follows immediately and must be met, because it threatens to dissolve the fourth step altogether. If reproduction thickens only insofar as the prior three feed it, if its health is a function of whether perception, reception, and sharing are occurring, then reproduction seems to be nothing in its own right, a mere name for the settling of what the other three produce, and not a distinct movement at all. The model, on this objection, has not four steps but three, with reproduction a redundant label for their result. The objection would hold if reproduction were merely the passive deposit of the prior three, but it is not, and the distinction is precise. Reproduction is the distinct movement by which what the prior three produce is allowed to settle as a raised starting point rather than dissipating, and this allowing is itself an act, distinct from the perceiving and receiving and sharing whose product it settles. The prior three can occur and their product can still fail to settle, dissipating instead of raising the next beginning, and what makes the difference is reproduction, the distinct movement of letting the produced beauty thicken into the field rather than pass and be lost. A couple may perceive, receive, and share a beauty and then let it dissipate, treat it as a passing pleasure complete in itself, and in that case the prior three occurred and reproduction did not, the beauty had and not settled. Reproduction is the distinct act of settling, the turning of a passing event into a raised ground, and it is no more redundant for depending on the prior three than a harvest is redundant for depending on the growing, the harvest being the distinct act without which the growing yields nothing kept. The fourth step depends on the first three and is not reducible to them, because the settling it performs is an act the others do not contain, the act by which their product is made to thicken rather than pass.

This deepens rather than dissolves the closing of the cycle, because it shows reproduction to be both dependent and distinct, the movement that both requires the prior three and adds to them the settling without which they would yield no accumulation. The cycle closes through reproduction not because reproduction passively receives the prior three but because reproduction actively settles their product into the raised ground that the next perception begins from, and this active settling is the hinge on which the line becomes a cycle, the point at which what was produced is turned back into a condition for what will be produced. Without the fourth step the first three would be a line, producing beauties that dissipate; with it they are a cycle, producing beauties that settle into a ground that raises the next production. Reproduction is thus the distinct movement that makes the difference between a line and a cycle, dependent on the prior three for its material and distinct from them in the settling it performs, and the model has four steps and not three because the settling is an act and not a mere result.

With this the figure that has organised the whole account completes itself. Each of the four sections met a conflict that threatened the model, and each conflict inverted on the same distinction, between a grasp that possesses and a generation that does not: non-division grounding discrimination, responding-yet-not-storing grounding accumulation, the preservation of difference grounding union, and now the feeding of reproduction by the living prior steps grounding the cycle's gain against the wheel. The four are not four lessons but one, the ancient counsel to generate without possessing and to bring to maturity without lording over, traced through perception, welcome, the between, and time. And the fourth completes the others, because it is the one that shows the model to be a cycle, returning to its root and not its origin, so that the form of the four steps is the very form of the good cycle they were meant to cultivate. The account, at its close, enacts what it describes; the paper, like the poetry the series took as its first paradigm, is built in the shape of its own claim.

\subsection{A layered synthesis}
\label{sub:reproduction-synthesis}

The accounts of reproduction take their places without identification. The philosophical analyses of repetition are taken at the level of metaphysics: they give the standing of the distinction between the spiral and the wheel, the forward repetition that gains against the backward one that retrieves and the differential repetition against the mechanical same, and they may be read alongside the holonomy criterion without being collapsed into it. The political economy of the reproductive act is taken at the level of conditions and justice: it gives the cause of degeneration that lies outside the laborer, the capture of the cycle's yield by an external power, and it guards the account against preaching inwardness to the dispossessed. Sedimentation and daily renewal are taken at the level of temporal structure and of cultivation: they give the form of an accumulation that does not hoard and a repetition that does not dull, and they resolve the standing worry that accumulation must violate non-abiding. None is written as equal to another, because a metaphysics of repetition, an analysis of ownership, and a discipline of daily renewal are not one thing, and the framework asks of each only what it can give.

What the synthesis secures is the close of the cycle and so of the model. Reproduction, rightly understood, is neither the mechanical making-again that the word first suggests nor a solitary feat of accumulation, but the settling, into a raised and non-hoarded ground, of what the living cycle has freshly produced, kept from alienation by the return of its yield to the coupling and kept from dulling by the continued turning of the prior three steps. The cultivation of the field, begun in perception and carried through reception and sharing, completes itself here not by reaching an end but by raising a beginning, and the four movements stand revealed as the smallest form in which a good cycle is enacted between two people, a form that gains by returning to its root and never to its origin.

\subsection{The reach and the defeat of each account}
\label{sub:reproduction-ledger}

The accounts of reproduction differ in whether they grasp its metaphysics, its political conditions, or its temporal form, and the accounting of their reach completes the section's argument.

The philosophical analyses of repetition explain the metaphysics of the distinction the section needs, the difference between the repetition that gains and the repetition that recurs, and they give the spiral and the wheel their standing in the history of thought, the forward repetition that renews and the differential repetition that carries difference against the mechanical same. What defeats them, taken alone, is their abstraction from the conditions under which repetition is actually one or the other. They tell us what good and degenerate repetition are but not what makes a given repetition degenerate in a given life, and pressed for that they would locate the difference in the inwardness of the repeater, which is exactly the sentimental error the political economy of the reproductive act exists to refuse. Their strength is the metaphysics and their defeat is the silence about conditions that the political-economic account must break.

The political economy of the reproductive act explains what the metaphysics omits, the condition that makes repetition deaden, the capture of its yield by an external power, and it guards the account against preaching inwardness to the dispossessed. What defeats it, taken alone, is its silence about the metaphysics it presupposes. It tells us that captured repetition deadens but draws on the distinction between the gaining and the recurring repetition that only the metaphysical account supplies, and it cannot by itself say what the accumulation is that capture diverts, needing the holonomy that the framework provides. Its strength is the condition and its defeat is the dependence on a metaphysics it does not itself furnish.

The accounts of sedimentation and daily renewal explain the temporal form of the accumulation, how it settles without hoarding and freshens without dulling, and they resolve the standing worry that accumulation must violate non-abiding. What defeats them, taken alone, is their silence about both the metaphysics and the conditions, describing the form of good temporal accumulation without grounding the distinction it instances or naming the conditions under which it fails. Their strength is the temporal form and their defeat is their dependence on the metaphysics and the political economy that place the form they describe.

The materialist resolution holds that reproduction is a real temporal process in a real field, possessing at once a metaphysical structure, a political condition, and a temporal form, and that the accumulation it secures is a real material accumulation whose ownership is a real material question. It is materialist in taking the accumulation to be real, a real change in the curvature of a real field, and the capture of its yield to be a real material relation of ownership, not a metaphor; it is dialectical in holding the metaphysics, the condition, and the form to be aspects of one real process whose truth is their articulation. Against an idealism that would make repetition's renewal a matter of the repeater's attitude alone, it holds the material condition of the yield's ownership; against a reduction that would make reproduction mere biological or economic recurrence, it holds the metaphysical distinction of the gaining from the recurring; against a formalism that would describe the temporal pattern without grounding it, it holds the metaphysics and the political economy that place the form. Reproduction is the real accumulation of a real field across time, structured and conditioned and formed, and reducible to none of these, which is the materialist dialectic applied to repetition without return. With this the four battlefields are complete, and the same materialist refusal of any theory's final authority has governed all four, each resolved not by a synthesis that crowns one account but by the determinate articulation of aspects, each doing its real work and none permitted to be the whole.
